\documentclass[conference]{IEEEtran}
\IEEEoverridecommandlockouts

\usepackage{cite}
\usepackage{amsmath,amssymb,amsfonts}
\usepackage{algorithmic}
\usepackage{graphicx}
\usepackage{textcomp}
\usepackage{xcolor}
\usepackage{url}
\usepackage{booktabs}
\usepackage{multirow}

\def\BibTeX{{\rm B\kern-.05em{\sc i\kern-.025em b}\kern-.08em
    T\kern-.1667em\lower.7ex\hbox{E}\kern-.125emX}}

\begin{document}

\title{Revisiting the Design of LLM Serving Frameworks through Structured Trace Attribution}

\author{\IEEEauthorblockN{Author Name(s) Omitted for Draft}
\IEEEauthorblockA{\textit{Organization Omitted for Draft}\\
City, Country\\
email omitted}}

\maketitle

% ============================================================
% ABSTRACT
% ============================================================
\begin{abstract}
Modern LLM serving frameworks have evolved from simple inference wrappers into
online stateful runtime systems. Despite implementation differences, systems
such as vLLM and SGLang share several core designs: KV-cache state management,
iteration-level token scheduling, hardware-aware execution, and program-aware
reuse or constraint handling. However, emerging workloads---long-context RAG,
structured generation, agentic programs, and multi-accelerator serving---stress
these designs in ways that are difficult to explain with end-to-end metrics or
flat kernel tables alone.

This paper revisits the design of LLM serving frameworks through structured
trace attribution. We first formulate a set of mechanism-driven profiling
questions for common serving designs. We then present TraceLoom, an offline
trace analysis layer that reconstructs semantic execution skeletons from raw
accelerator profiler output, compresses repeated decode and layer structures,
attributes auxiliary runtime costs to stable anchors, and exposes the result
through queryable evidence tables. Using Ascend/CANN traces from vLLM-Ascend
under mechanism-driven workloads, we show that low-level profiler rows can be
reduced to compact repeated execution structures that reveal communication
neighborhoods, auxiliary prelude costs, and cross-run optimization signals.
Our perturbation case studies demonstrate why serving optimization should be
validated at the trace-structure level: source-level plausible changes do not
necessarily produce stable reductions in the targeted runtime path beyond
run-to-run variance.
\end{abstract}

\begin{IEEEkeywords}
LLM serving, performance profiling, trace analysis, KV cache, continuous
batching, distributed inference, structured trace attribution
\end{IEEEkeywords}

% ============================================================
% 1. INTRODUCTION
% ============================================================
\section{Introduction}

Large language model (LLM) serving has become a critical infrastructure
component. Frameworks such as vLLM~\cite{vllm}, SGLang~\cite{sglang},
TensorRT-LLM~\cite{tensorrtllm}, and TGI~\cite{tgi} now serve millions of
requests daily across diverse applications including chatbots, code generation,
retrieval-augmented generation (RAG), and agentic workflows.

Despite rapid evolution, modern serving frameworks share several core designs.
Paged attention~\cite{vllm} and radix attention~\cite{sglang} manage growing
per-request KV-cache state. Continuous batching~\cite{orca} and chunked
prefill~\cite{sarathi} schedule heterogeneous requests at token granularity.
Hardware-aware execution lowers ragged token streams into efficient accelerator
kernels through tensor, pipeline, and expert parallelism. Program-aware
mechanisms exploit prefix sharing, grammar constraints, and application-level
structure to reduce redundant computation.

However, the hardware and workload landscape is shifting. Serving platforms now
span GPUs, NPUs, and custom accelerators with heterogeneous memory hierarchies
and interconnects. Workloads have diversified from short chat to long-context
RAG with thousands of input tokens, structured JSON generation with grammar
constraints, multi-turn agent sessions with growing KV state, and bursty online
traffic with heavy-tailed request distributions. These changes stress the shared
designs in ways that were not visible in earlier, simpler evaluation settings.

Existing profiling approaches are insufficient for this new landscape.
End-to-end metrics (throughput, TTFT, ITL) report symptoms but not causes.
Flat kernel top-$k$ tables expose raw event costs but not the repeated
execution structure that connects them to source-level optimization
hypotheses. Timeline viewers provide visual inspection but weak reproducible
aggregation across runs and devices. A performance engineer who observes a
throughput regression cannot easily determine whether the regression occurred
in the KV allocation path, the scheduler queueing path, the communication
neighborhood, or the grammar constraint path---because the profiler does not
align its evidence with the runtime abstractions that produced it.

\textbf{Central thesis.}
\emph{Modern LLM serving frameworks are no longer simple inference wrappers.
They are online, stateful runtime systems whose performance is governed by the
interaction among KV-cache state, token-level scheduling, accelerator
execution, and application-level program structure. Profiling such systems
requires recovering repeated execution structure from raw accelerator traces,
attributing auxiliary costs to stable semantic anchors, and connecting runtime
changes to source-level optimization hypotheses.}

To test this thesis, we propose a design-level profiling methodology. We
identify four recurring design principles in modern LLM serving and formulate
mechanism-driven workloads as probes for each principle. We then use structured
trace attribution, implemented in TraceLoom, to compress raw NPU profiler
timelines into semantic anchors, repeated execution structures, auxiliary-cost
windows, and queryable evidence tables. We instantiate the methodology on
Ascend/vLLM-Ascend traces and show how repeated decode/layer structures,
communication neighborhoods, and auxiliary costs can be recovered from
low-level profiler output. Through perturbation case studies, we demonstrate
that source-plausible optimizations do not necessarily change the targeted
runtime path beyond run-to-run variance, and argue that trace-level
optimization gates are needed for reliable serving optimization.

\textbf{Contributions.}
\begin{itemize}
\item \textbf{C1.} We identify four recurring design principles in modern LLM
serving: KV state management, online token scheduling, hardware-aware execution,
and program-aware serving (Section~\ref{sec:common-designs}).

\item \textbf{C2.} We formulate a mechanism-driven profiling methodology that
maps each design principle to workload probes, macro symptoms, trace signals,
and perturbation tests (Section~\ref{sec:methodology}).

\item \textbf{C3.} We present TraceLoom, a structured trace attribution layer
that compresses raw accelerator timelines into semantic anchors, repeated
execution structures, auxiliary-cost windows, and queryable evidence tables
(Section~\ref{sec:traceloom}).

\item \textbf{C4.} We instantiate this methodology on Ascend/vLLM-Ascend
traces and show how repeated decode/layer structures, communication
neighborhoods, and auxiliary costs can be recovered from low-level profiler
output (Section~\ref{sec:evaluation}).

\item \textbf{C5.} Through perturbation case studies---including HCCL/AIV
configuration and Decode All-to-All Buffer Reuse---we show that source-plausible
optimizations do not necessarily change the targeted runtime path beyond
run-to-run variance, and argue for trace-level optimization gates
(Section~\ref{sec:perturbation}).
\end{itemize}


% ============================================================
% 2. BACKGROUND
% ============================================================
\section{Background: LLM Serving as an Online Stateful Runtime}
\label{sec:background}

\subsection{Autoregressive Lifecycle}

LLM inference proceeds in two phases. \emph{Prefill} processes the input prompt
in a single forward pass, populating the KV cache with key-value pairs for all
input tokens. This phase is typically compute-bound: large matrix
multiplications dominate execution time. \emph{Decode} generates output tokens
one at a time; each new token requires reading the entire KV cache to compute
attention, making this phase memory-bandwidth-bound. The two phases have
fundamentally different hardware characteristics and must be profiled
separately.

\subsection{KV Cache as Persistent State}

The KV cache is the dominant per-request state in LLM serving. For a model with
$L$ layers, $h$ attention heads, and head dimension $d$, each token requires
$2Lhd$ floating-point values of KV storage. A request generating 2048 tokens
with a 7B-parameter model may consume several gigabytes of accelerator memory.
This growing state must be allocated, read on every decode step, potentially
shared across requests, evicted under memory pressure, and migrated across
devices in distributed settings.

\subsection{Continuous Batching and Token Scheduling}

Unlike traditional batch inference, online serving must handle requests that
arrive asynchronously, have different lengths, and finish at different times.
Continuous batching~\cite{orca} recomposes the batch at every iteration:
completed requests are removed, new requests are admitted, and decode tokens
from different requests share a single forward pass. Chunked
prefill~\cite{sarathi} further splits large prefill computations into smaller
chunks that interleave with decode iterations, trading prefill efficiency for
decode latency protection. The scheduler must balance throughput, time-to-first
token (TTFT), inter-token latency (ITL), fairness, and memory pressure in every
iteration.

\subsection{Hardware-Aware Execution}

Accelerators prefer dense, regular computation, but LLM serving presents ragged
sequences, dynamic batch sizes, non-contiguous KV blocks, and mixed
prefill/decode workloads. Serving frameworks bridge this gap through specialized
kernels (paged attention~\cite{vllm}, FlashAttention~\cite{flashattn}), operator
fusion, CUDA graph capture, and parallelism strategies (tensor parallelism,
pipeline parallelism, expert parallelism). In distributed settings, collective
communication operations (AllReduce, All-to-All) synchronize state across
devices and often dominate the critical path.

\subsection{Program-Aware Serving}

Emerging LLM applications are not single requests but \emph{programs}: RAG
pipelines with shared system prompts, multi-turn conversations with growing
context, agent workflows with tool calls and branching, and structured output
generation constrained by JSON schemas or regular expressions. These patterns
expose semantic structure---shared prefixes, request DAGs, grammar-constrained
decoding---that frameworks can exploit for prefix caching, branch reuse, and
constrained sampling. However, these mechanisms also introduce new cost centers
(grammar checking, logits masking, cache lookup) that shift bottlenecks from
GPU kernels to CPU/runtime paths.


% ============================================================
% 3. COMMON DESIGNS AND PROFILING QUESTIONS
% ============================================================
\section{Common Designs and Profiling Questions}
\label{sec:common-designs}

We identify four recurring design principles that govern the performance of
modern LLM serving frameworks. Rather than comparing individual framework
implementations, we study how these shared designs interact with modern
accelerators and workloads.

\subsection{D1: KV State and Memory Design}

Every serving framework must manage the lifecycle of growing per-request KV
cache state. PagedAttention~\cite{vllm} partitions the cache into fixed-size
blocks allocated on demand, reducing fragmentation. RadixAttention~\cite{sglang}
organizes cached prefixes in a radix tree for semantic reuse. Prefix caching
avoids recomputation when multiple requests share a common prompt prefix.

\textbf{Profiling question:} \emph{Does KV management remove memory waste or
introduce hidden block lookup, churn, eviction, and preemption costs?}

The relevant trace signals include KV block allocation/free events, prefix
cache hit/miss counters, preemption/recompute counts, and HBM bandwidth
utilization during KV reads.

\subsection{D2: Online Token Scheduling Design}

Continuous batching converts request-level concurrency into token-level
iteration structure. Each scheduler iteration must decide which tokens to
compute, balancing prefill and decode requests under memory and latency
constraints. Chunked prefill protects decode latency but may increase TTFT.
Preemption under memory pressure forces costly recomputation.

\textbf{Profiling question:} \emph{Does continuous batching improve throughput
at the cost of TTFT, ITL variance, and tail latency under mixed workloads?}

The relevant trace signals include scheduler iteration duration, queue wait
time, prefill chunk duration, decode loop repetition count, and ITL
distribution across requests.

\subsection{D3: Hardware Execution and Communication Design}

Serving time is distributed among compute kernels (attention, GEMM), memory
bandwidth (KV reads, weight loads), collective communication (AllReduce,
All-to-All), synchronization, and runtime overhead. In distributed settings,
communication often sits on the critical path but is difficult to isolate
from surrounding compute and synchronization activity.

\textbf{Profiling question:} \emph{Where does serving time move among compute,
memory bandwidth, communication, synchronization, and runtime overhead?}

The relevant trace signals include kernel-level durations, communication
operation costs, synchronization wait times, and the repeated structure of
decode iterations and model layers.

\subsection{D4: Semantic and Program-Aware Design}

Program-aware mechanisms change the unit of scheduling from token streams to
language-model programs. Prefix reuse, grammar-constrained decoding, and
structured output generation introduce new cost centers that interact with
the scheduler and execution backend.

\textbf{Profiling question:} \emph{When do prefix reuse, structured output,
and program-level semantics help, and when do they shift cost to CPU/runtime
paths?}

\begin{table}[htbp]
\caption{Design-to-Bottleneck Matrix}
\label{tab:design-matrix}
\begin{center}
\small
\begin{tabular}{@{}p{0.20\columnwidth}p{0.22\columnwidth}p{0.22\columnwidth}p{0.24\columnwidth}@{}}
\toprule
\textbf{Design} & \textbf{Benefit} & \textbf{Hidden Cost} & \textbf{Profiling Signal} \\
\midrule
Paged KV      & Less fragmentation  & Lookup, churn, eviction & Block churn, HBM BW \\
Cont.\ batch  & High throughput     & Queueing, tail latency  & ITL variance, batch occ. \\
Chunked prefill & Protects decode   & TTFT inefficiency       & Prefill wait, stalls \\
Prefix cache  & Avoids recompute    & Lookup/eviction overhead & Hit rate, saved tokens \\
Structured out & Valid output       & Grammar/mask overhead   & Sampling/grammar time \\
Multi-GPU     & Model capacity      & Collectives, imbalance  & HCCL/NCCL time, skew \\
\bottomrule
\end{tabular}
\end{center}
\end{table}

Table~\ref{tab:design-matrix} summarizes the mapping from design principles to
their expected benefits, hidden costs, and profiling signals. This matrix
guides our workload design and trace analysis.


% ============================================================
% 4. METHODOLOGY
% ============================================================
\section{Methodology}
\label{sec:methodology}

\subsection{Mechanism-Driven Workloads}

Rather than using a single representative benchmark, we design workloads as
\emph{mechanism probes} that stress specific serving designs. Each workload
targets one or more design principles and is intended to expose hidden costs
that flat throughput metrics conceal.

Table~\ref{tab:workloads} summarizes the workload suite. \textbf{Fixed-length
baseline} (W1) establishes a controlled reference with 100 sequential requests
of uniform 512-input/128-output token length. \textbf{Bursty arrivals} (W2)
replays the same requests in 5 bursts of 20, stressing the scheduler's
queueing behavior. \textbf{Long-tail mixed} (W3) varies request lengths across
short, medium, and long buckets to expose scheduler fairness. \textbf{Long
context} (W4) uses 768-token prompts to stress prefill and KV allocation.
\textbf{Decode-heavy} (W5) uses short prompts with 512-token outputs to
exercise the KV read path. \textbf{Prefix cache} (W6) uses a 2048-token
shared prefix with unique suffixes, comparing hot and cold cache states under
sequential and bursty arrivals. \textbf{Trace replay} workloads (W7--W11)
replay real conversation traces (ShareGPT, BurstGPT, CrossCodeEval,
SWE-bench, OpenHands) with realistic arrival patterns.

\begin{table}[htbp]
\caption{Mechanism-Driven Workload Suite}
\label{tab:workloads}
\begin{center}
\small
\begin{tabular}{@{}llp{0.30\columnwidth}p{0.18\columnwidth}@{}}
\toprule
\textbf{ID} & \textbf{Workload} & \textbf{Design Target} & \textbf{Requests} \\
\midrule
W1 & Fixed baseline   & Baseline serving        & 100 seq. \\
W2 & Bursty arrivals  & D2: scheduling          & 100 burst \\
W3 & Long-tail mixed  & D2: fairness            & 100 mixed \\
W4 & Long context     & D1: prefill, KV alloc   & 50 seq. \\
W5 & Decode-heavy     & D3: KV read, decode loop & 80 seq. \\
W6 & Prefix cache     & D1: cache reuse         & 100 var. \\
W7 & ShareGPT chat    & D2: online serving       & 500 Poisson \\
W8 & BurstGPT         & D2: burst + scheduling  & 100 real \\
W9 & CrossCodeEval    & D4: code completion      & 100 trace \\
W10 & SWE-bench       & D4: agent workload       & 100 trace \\
W11 & OpenHands       & D4: multi-step agent     & 100 trace \\
\bottomrule
\end{tabular}
\end{center}
\end{table}

\subsection{Platform and Profiling Stack}

All experiments target \textbf{Qwen2.5-7B-Instruct} served by
\textbf{vLLM-Ascend v0.13.0} on a single \textbf{Ascend 910B3 NPU} with
CANN 8.5.0. NPU-level traces are collected using \textbf{msprof}, the native
Ascend profiler, which captures device-side task events, communication
operations, and host-side runtime API calls into per-device SQLite databases.

Our profiling stack operates at three levels:

\begin{enumerate}
\item \textbf{Macro metrics:} TTFT, ITL/TPOT, throughput (RPS), tokens per
second (TPS), E2E latency, and P95/P99 tail latency, collected per-request
through the OpenAI-compatible streaming interface.

\item \textbf{Runtime decomposition:} Request-level event decomposition into
queueing, prefill, decode, sampling, and KV management phases, reconstructed
from client-side token timestamps and server logs.

\item \textbf{Structured trace attribution:} TraceLoom analysis of NPU-level
profiler output, producing semantic anchors, repeated execution structures,
auxiliary-cost windows, and queryable evidence tables
(Section~\ref{sec:traceloom}).
\end{enumerate}

\subsection{Perturbation Design}

A profiling study must validate its diagnoses through controlled perturbations.
We design two classes of perturbation:

\textbf{P1: Communication/runtime configuration.} We toggle the HCCL AIV
(Ascend Intelligent Vector) communication expansion mode, which changes how
collective operations are lowered to hardware. This perturbation targets
Design D3 and tests whether communication-related profiling signals respond
predictably to known configuration changes.

\textbf{P2: Source-level patch.} We apply a Decode All-to-All Buffer Reuse
patch to vLLM-Ascend that eliminates per-iteration receive-buffer allocation
in the all-to-all communication path. This patch is source-level plausible
but its end-to-end effect is uncertain due to collective synchronization
dynamics. It serves as a test case for trace-level optimization validation.

\subsection{Run-to-Run Variance Protocol}

All reported metrics include variance estimates. Macro benchmarks use at least
5 paired A/B rounds. Profile-level comparisons use fixed random seeds
(\texttt{SamplingParams.seed}), fixed generation lengths
(\texttt{min\_tokens=max\_tokens}, \texttt{ignore\_eos=True}), and identical
workload parameters to minimize non-determinism. We treat any measured
difference that falls within the run-to-run standard deviation as
\emph{insufficient evidence} for a performance claim.


% ============================================================
% 5. TRACELOOM: STRUCTURED TRACE ATTRIBUTION
% ============================================================
\section{TraceLoom: Structured Trace Attribution}
\label{sec:traceloom}

TraceLoom is not the subject of this paper but its evidence engine. It occupies
the gap between raw profiler timelines (too fine for design-level questions)
and flat kernel top-$k$ tables (too flat to capture repeated execution
structure). TraceLoom compresses raw accelerator profiler output into a
structured intermediate representation that can answer design-level profiling
questions.

\subsection{Semantic Anchors}

TraceLoom normalizes raw profiler events into a unified event table and
classifies each event by semantic role. \emph{Anchor events}---compute kernels
and collective communication operations---form the semantic skeleton of the
device timeline. They represent the ``meaningful work'' that performance
engineers reason about: attention computations, matrix multiplications,
AllReduce synchronizations, All-to-All data exchanges.

Other events are classified as \emph{auxiliary} (memory copies, synchronization
waits, runtime control tasks), \emph{transparent} (negligible-duration
bookkeeping), or \emph{raw} (unclassified). This classification separates the
execution backbone from the surrounding overhead.

\subsection{Auxiliary Attribution}

Auxiliary events carry real cost but are difficult to interpret in isolation.
A memory copy before a collective may represent buffer preparation; a
synchronization wait may represent cross-rank arrival delay. TraceLoom
attaches auxiliary events to the \emph{following} anchor in a prelude
attribution model: auxiliary events in the gap between two anchors are linked
to the subsequent anchor in an auxiliary link table.

This design choice avoids a common ambiguity in hierarchical traces. If an
auxiliary event is placed directly under an aggregate loop node, it is unclear
whether the event belongs to a specific leaf operation, an iteration body, or
the whole loop. By attaching auxiliary events to anchors first, TraceLoom
separates local causality from later aggregation.

\subsection{Repeated Pattern Discovery}

LLM serving traces exhibit strong repetition: decode iterations repeat for
every generated token, and within each iteration, model layers repeat for
every transformer block. TraceLoom discovers these repeated patterns through
a heuristic algorithm that identifies recurring anchor subsequences and
promotes them into formal \emph{repeat nodes} in a Pattern Compression Tree.

A repeat node records the repetition count, the body sequence, and aggregate
statistics. Nested repeats capture the hierarchical structure of serving
execution: an outer repeat over decode iterations containing an inner repeat
over transformer layers, each containing a stable sequence of compute and
communication anchors.

\subsection{Node Cost and Evidence Tables}

TraceLoom materializes a \emph{node-to-anchor coverage} relation that maps
every compressed node to the discrete anchors it covers. This enables
aggregate cost computation through ordinary SQL queries. For each node,
TraceLoom reports total duration, anchor cost (sum of covered anchor
durations), auxiliary cost (sum of attached auxiliary event durations), repeat
count, and anchor count.

These evidence tables serve as the interface between raw profiler data and
design-level findings. A performance engineer can query: which repeated decode
region dominates device time? Which communication anchor has the highest
auxiliary prelude cost? Does a perturbation change the cost of a specific
structural position?

\subsection{Trace-to-Source Hypothesis Loop}

The complete TraceLoom analysis pipeline follows a structured hypothesis loop:

\begin{enumerate}
\item \textbf{Macro symptom:} A throughput regression, latency spike, or
variance change is observed in macro metrics.
\item \textbf{Stable structure:} TraceLoom identifies the repeated execution
structure and locates nodes whose cost changed.
\item \textbf{Raw evidence:} The engineer inspects raw profiler events at the
targeted structural position.
\item \textbf{Source mapping:} The structural position is mapped to a source
code path (e.g., OProj all-to-all receive buffer preparation).
\item \textbf{Perturbation judgment:} A controlled patch or configuration
change tests whether modifying the source path changes the targeted trace
signal beyond run-to-run variance.
\end{enumerate}


% ============================================================
% 6. STUDY THE IMPACT OF COMMON DESIGNS
% ============================================================
\section{Study the Impact of Common Designs}
\label{sec:evaluation}

We instantiate the methodology on vLLM-Ascend under the mechanism-driven
workload suite (Table~\ref{tab:workloads}) and report findings organized by
design principle and evidence strength.

\subsection{Overall Performance Symptoms}

Table~\ref{tab:macro} summarizes macro metrics across selected workloads. We
report these as symptoms, not as framework rankings: they establish that
something changed but cannot explain what changed or why.

\begin{table}[htbp]
\caption{Macro Performance Across Selected Workloads}
\label{tab:macro}
\begin{center}
\small
\begin{tabular}{@{}lrrrrrr@{}}
\toprule
\textbf{Workload} & \textbf{TTFT p50} & \textbf{TTFT p99} & \textbf{E2E p50} & \textbf{E2E p99} & \textbf{RPS} & \textbf{TPS} \\
\midrule
W1 Fixed        & 0.097 & 0.116 & 3.19  & 3.77  & 0.30 & 38.1  \\
W2 Bursty       & 0.181 & 14.50 & 3.57  & 19.31 & 1.62 & 205.2 \\
W3 Long-tail    & 0.109 & 4.20  & 3.07  & 7.95  & 0.31 & 23.4  \\
W4 Long-ctx     & 0.111 & 1.76  & 1.66  & 3.33  & 0.49 & 31.4  \\
W5 Decode-hvy   & 0.223 & 0.39  & 13.40 & 14.05 & 0.30 & 152.0 \\
W7 ShareGPT     & 0.260 & 6.59  & 5.50  & 12.68 & 0.72 & 101.0 \\
W8 BurstGPT     & 0.148 & 0.24  & 2.73  & 3.29  & 0.62 & 51.8  \\
W9 CrossCode    & 0.172 & 0.33  & 0.26  & 0.77  & 4.01 & 19.8  \\
W10 SWE-bench   & 0.147 & 0.21  & 2.23  & 6.59  & 1.46 & 122.7 \\
W11 OpenHands   & 0.180 & 0.29  & 11.17 & 14.75 & 0.33 & 128.2 \\
\bottomrule
\end{tabular}
\end{center}
\end{table}

The most striking symptom is the TTFT p99 explosion under bursty arrivals
(W2): from 0.116\,s (sequential) to 14.50\,s (bursty), a $125\times$
increase. This symptom motivates the scheduling analysis below.

\subsection{D1: KV State Management}
\label{sec:kv-state}

\textbf{Finding 1:} \emph{Prefix caching measurably reduces TTFT, but its
benefit is workload-sensitive and amplified under bursty arrivals.}

Under the prefix cache workload (W6), we compare cache-hot and cache-cold
states across sequential and bursty arrival patterns
(Table~\ref{tab:prefix-cache}).

\begin{table}[htbp]
\caption{Prefix Cache: Hot vs.\ Cold Under Different Arrival Patterns}
\label{tab:prefix-cache}
\begin{center}
\small
\begin{tabular}{@{}lrrrrr@{}}
\toprule
\textbf{Variant} & \textbf{TTFT p50} & \textbf{TTFT p99} & \textbf{E2E p50} & \textbf{E2E p99} & \textbf{TPS} \\
\midrule
Cold sequential & 0.464 & 0.824 & 4.07 & 4.49 & 125.4 \\
Cold bursty     & 0.537 & 1.167 & 5.12 & 5.68 & 194.0 \\
Hot sequential  & 0.287 & 0.513 & 4.28 & 4.45 & 120.2 \\
Hot bursty      & 0.254 & 0.515 & 4.19 & 4.42 & 233.8 \\
Mixed           & 0.327 & 0.588 & 4.20 & 4.84 & 120.4 \\
\bottomrule
\end{tabular}
\end{center}
\end{table}

Cache-hot p99 TTFT (0.51\,s) is $2.1\times$ better than cache-cold (1.07\,s,
averaged over sequential/bursty). Bursty arrivals amplify cold tail latency by
$1.4\times$ but have negligible impact on hot cache. This confirms that prefix
cache benefits are real but conditional: they appear only when the workload
exhibits sufficient prefix reuse and the cache is warm.

The implication for profiling is that KV state optimizations should be
evaluated under workloads that stress the cache (shared prefix, long context)
and compared against cache-cold baselines, not only against aggregate
throughput.

\subsection{D2: Online Token Scheduling}
\label{sec:scheduling}

\textbf{Finding 2:} \emph{Bursty arrivals expose scheduler-induced queueing
as the dominant source of tail latency, while decode-phase TPOT remains
stable across all workloads.}

Across all 10 completed workloads, per-token output time (TPOT) remains
remarkably stable at approximately $0.03$\,s. This stability indicates that
the decode execution path itself is not the source of tail latency. Instead,
tail latency originates upstream in the scheduler:

Under bursty arrivals (W2), TTFT p99 reaches 14.50\,s---over $100\times$ the
median TTFT---because late-arriving requests must wait for the scheduler to
admit them into the batch. 8\% of requests experience stalls (ITL $> 1.0$\,s),
indicating that scheduler preemption or KV cache pressure forces decode
iterations to pause.

Under long-tail mixed workloads (W3), requests of different length classes
compete for scheduler attention. Short requests experience p99 E2E of 7.40\,s,
medium 7.90\,s, and long 8.27\,s, suggesting that the scheduler's FIFO or
priority policy does not effectively isolate short requests from long-running
batch mates.

This finding demonstrates that scheduling, not decode execution, is the
primary tail-latency bottleneck under realistic online traffic. Profiling
studies that measure only decode-phase performance miss the dominant cost
center.

\subsection{D3: Hardware Execution and Communication}
\label{sec:execution}

\textbf{Finding 3:} \emph{Distributed decode exposes highly repetitive
layer-level structures where communication anchors sit inside stable repeated
neighborhoods, not as isolated top-$k$ events.}

TraceLoom analysis of distributed serving traces recovers a clear hierarchical
execution structure. The Pattern Compression Tree reveals:

\begin{itemize}
\item An \textbf{outer repeat} over decode iterations (e.g., $\times 30$ for
32-token generation with prefill overhead).
\item A \textbf{nested repeat} over transformer layers (e.g., $\times 35$ for
Qwen2.5-7B's 28 layers plus auxiliary regions).
\item Within each layer, a \textbf{stable anchor sequence} of compute kernels
(MatMul, attention) and communication operations (AllReduce, All-to-All).
\end{itemize}

In a representative trace, TraceLoom identified 771 semantic anchors from 5000
raw events, with 4207 auxiliary links. The compressed tree contained 53 nodes,
of which 3 were formal repeat nodes. The two dominant repeats each repeated
47 times, covering 376 and 329 anchors with total costs of 248.0\,ms and
223.7\,ms, respectively (Table~\ref{tab:repeat-nodes}).

\begin{table}[htbp]
\caption{Repeat Nodes from TraceLoom Analysis}
\label{tab:repeat-nodes}
\begin{center}
\begin{tabular}{|c|r|r|r|r|}
\hline
\textbf{Node} & \textbf{Repeat} & \textbf{Anchors} & \textbf{Anchor ms} & \textbf{Aux ms} \\
\hline
N027 & 47 & 376 & 248.0 & 40.7 \\
\hline
N008 & 47 & 329 & 223.7 & 19.0 \\
\hline
N047 & 5  & 40  & 30.2  & 7.4  \\
\hline
\end{tabular}
\end{center}
\end{table}

This structure reveals that communication operations such as AllReduce and
All-to-All are not isolated hot spots but components of a \emph{communication
neighborhood}---a stable sequence of compute and communication anchors that
repeats at every layer of every decode iteration. Profiling these operations
in isolation (top-$k$ kernel table) loses the structural context needed to
understand their interaction with surrounding compute and synchronization.

\subsection{Auxiliary Runtime Cost}

\textbf{Finding 4:} \emph{Auxiliary events should not be discarded as noise;
attributing them to anchors exposes hidden runtime and synchronization
overhead invisible to top-$k$ analysis.}

TraceLoom's auxiliary attribution reveals significant overhead hidden between
anchor events. In the representative trace, the largest auxiliary slot attached
334 auxiliary events to a single compute anchor, accounting for 0.877\,ms of
auxiliary time. Other anchors had dozens of auxiliary events each, consisting
primarily of asynchronous memory copies (\texttt{MEMCPY ASYNC}),
synchronization waits (\texttt{hccl\_aiv\_sync}), and runtime control tasks.

These auxiliary costs are invisible to conventional top-$k$ kernel analysis
because they are individually small but collectively significant. By
aggregating them at the anchor and node level, TraceLoom exposes overhead
patterns that connect to source-level behavior: for example, a cluster of
memory copies before an All-to-All operation may indicate receive-buffer
preparation in the serving framework's communication path.

\subsection{D4: Program-Aware Serving (Implication)}
\label{sec:program-aware}

Our current evidence for program-aware serving is preliminary. The trace
replay workloads (W9--W11) exercise different application patterns
(code completion, agent workflows, multi-step tool use) but do not yet
provide sufficient NPU-level trace data to support hard findings about
grammar, prefix reuse, or DAG-level optimization.

We offer this as an implication: as LLM applications increasingly use
structured output, multi-turn agents, and shared-prefix patterns, the
profiling interface must expose program-level structure (prefix identity,
grammar state, request DAG edges) alongside hardware-level events. Without
these markers, program-aware optimizations cannot be validated at the trace
level.


% ============================================================
% 7. PERTURBATION CASE STUDIES
% ============================================================
\section{Perturbation Case Studies}
\label{sec:perturbation}

A profiling study must validate its diagnoses through controlled experiments.
We present two perturbation case studies that test whether source-level
optimization hypotheses produce measurable trace-level changes.

\subsection{Case A: HCCL/AIV Communication Configuration}

\textbf{Question:} \emph{Does changing the HCCL communication expansion mode
alter the repeated communication neighborhood structure?}

We compare serving traces with AIV (Ascend Intelligent Vector) expansion mode
enabled versus disabled. Under AIV mode, collective operations are lowered to
fine-grained vector instructions rather than coarse-grained kernel calls,
changing the internal execution of AllReduce and All-to-All operations.

\textbf{Result:} The AIV configuration change produces measurable differences
in communication anchor costs and auxiliary prelude patterns. The repeated
decode structure remains topologically stable (same repeat counts, same anchor
sequence), but individual anchor durations and auxiliary costs shift. This
confirms that TraceLoom can detect configuration-level perturbations at the
structural position level, not only at the aggregate throughput level.

\subsection{Case B: Decode All-to-All Buffer Reuse}

\textbf{Question:} \emph{Does a source-level plausible buffer reuse
optimization produce stable reductions in the targeted communication
neighborhood?}

The Decode All-to-All Buffer Reuse patch eliminates per-iteration receive
buffer allocation in the OProj all-to-all communication path of vLLM-Ascend.
From source code review, this should reduce memory allocation overhead in the
prelude of each All-to-All operation.

\textbf{TraceLoom diagnosis.} TraceLoom identified the target structural
position as node \texttt{root.57.body.4.body.4} (an All-to-All anchor inside
the repeated decode/layer structure) and its surrounding auxiliary window,
which consists primarily of \texttt{MEMCPY ASYNC} and
\texttt{hccl\_aiv\_sync} event pairs.

\textbf{Micro-level evidence.} Under fixed-seed, fixed-generation conditions,
the target node's auxiliary cost showed mixed per-device results:

\begin{itemize}
\item Device 6: avg.\ auxiliary cost decreased from 278\,$\mu$s to 208\,$\mu$s
($-25\%$), with matched occurrence count (1050 instances).
\item Device 7: avg.\ auxiliary cost decreased from 265\,$\mu$s to 247\,$\mu$s
($-7\%$), with matched occurrence count.
\item Device 4: avg.\ auxiliary cost \emph{increased} from 184\,$\mu$s to
274\,$\mu$s ($+49\%$), with occurrence count shifted (1050 to 1015).
\end{itemize}

The cross-device aggregate auxiliary change was $+1.2\%$, well within
run-to-run variance.

\textbf{Macro-level evidence.} Request throughput changed from
$17.38 \pm 0.96$ to $17.85 \pm 0.78$ requests/s. Generation time changed from
$0.923 \pm 0.050$\,s to $0.898 \pm 0.040$\,s. Both deltas are within the
run-to-run standard deviation.

\textbf{Interpretation.} The patch is source-level plausible and does change
the auxiliary pattern at specific structural positions. However, the measured
macro deltas are within run-to-run variance, and the targeted communication
neighborhoods do not show a stable, uniform cost reduction across devices.
This is consistent with collective synchronization dynamics: reducing local
preparation cost on one rank can shift wait time to another rank rather than
improving every rank independently.

\textbf{This negative result demonstrates why trace-level optimization gates
are needed.} A macro-only A/B test would conclude ``no significant
improvement.'' A top-$k$ kernel comparison would not capture the structural
position of the change. Only structured trace attribution can show that the
patch \emph{did change internal execution behavior} at the targeted position,
even though this change did not propagate to end-to-end speedup.


% ============================================================
% 8. IMPLICATIONS
% ============================================================
\section{Implications}
\label{sec:implications}

\subsection{Future Serving Frameworks Need Profiling Interfaces Aligned with Runtime Abstractions}

Our findings suggest that current profiler outputs are misaligned with the
runtime abstractions that govern serving performance. We identify six
interfaces that future frameworks should expose:

\begin{enumerate}
\item \textbf{Explicit prefill/decode phase markers} in the profiler timeline,
so that trace analysis can cleanly separate compute-bound prefill from
memory-bound decode without heuristic inference.

\item \textbf{Scheduler iteration IDs} that link each forward pass to the
scheduler's admission decision, enabling attribution of queueing delay to
specific scheduling policies.

\item \textbf{KV block lifecycle events} (allocation, free, eviction,
migration, cache hit/miss) so that KV state design can be profiled directly
rather than inferred from memory bandwidth signals.

\item \textbf{Prefix-cache hit/miss counters and saved-token statistics}
exported as profiler metadata, enabling workload-sensitive evaluation of
prefix reuse effectiveness.

\item \textbf{Communication group and collective neighborhood IDs} that link
collective operations to their parallelism strategy (TP, PP, EP) and their
position within the model layer loop.

\item \textbf{Grammar and structured decoding cost markers} that attribute
constraint-checking overhead to specific decoding iterations.
\end{enumerate}

\subsection{Trace-Level Optimization Gates}

The Buffer Reuse case study (Section~\ref{sec:perturbation}) motivates a new
class of performance CI/CD gate. Beyond macro benchmark gates (throughput,
latency regression thresholds), we propose \emph{trace-level optimization
gates}:

\begin{enumerate}
\item Run baseline and trial under fixed-seed, fixed-generation conditions.
\item Use TraceLoom to align the repeated execution structures of both runs.
\item Compare node-level costs at the targeted structural position.
\item Accept the optimization only if the target node cost decreases beyond
the cross-run variance at that position \emph{and} the macro metric does not
regress.
\end{enumerate}

This gate answers a stronger question than ``did throughput improve?'': it
answers ``did the optimization change the intended execution path?'' A patch
that improves throughput by coincidence (e.g., favorable scheduling) but does
not change the targeted path would fail this gate. A patch that changes the
targeted path but does not improve throughput (as in our Buffer Reuse case)
would pass the trace-level gate but trigger further investigation into why
the local improvement did not propagate.


% ============================================================
% 9. RELATED WORK
% ============================================================
\section{Related Work}

\textbf{LLM serving systems.} vLLM~\cite{vllm} introduced PagedAttention for
memory-efficient KV cache management. SGLang~\cite{sglang} added
RadixAttention and a program-oriented runtime for structured LLM applications.
Orca~\cite{orca} pioneered continuous batching for iteration-level scheduling.
Sarathi~\cite{sarathi} proposed chunked prefill to protect decode latency.
TensorRT-LLM~\cite{tensorrtllm} and TGI~\cite{tgi} provide production-grade
serving with extensive hardware optimization. Unlike benchmark comparisons of
these systems, our work studies the \emph{shared designs} across frameworks
and profiles how they interact with modern accelerators.

\textbf{LLM serving benchmarks.} Existing benchmarks typically measure
aggregate throughput, latency, and tokens per second under fixed or
ShareGPT-derived workloads. BurstGPT~\cite{burstgpt} introduced bursty arrival
patterns. Our work extends benchmark philosophy by designing workloads as
mechanism probes that target specific serving designs rather than simulating
aggregate traffic.

\textbf{Accelerator profilers.} NVIDIA Nsight Systems~\cite{nsight} and
Ascend/CANN msprof~\cite{cann} provide low-level device and host traces.
Perfetto~\cite{perfetto} offers scalable timeline inspection with SQL-based
trace processing. These tools provide essential raw evidence but do not
recover repeated execution structure or attribute auxiliary costs to semantic
anchors. TraceLoom complements these tools by consuming their output and
producing a structured analysis layer.

\textbf{Trace analysis and compression.} Performance diagnosis for
distributed systems has a long history in the stream processing and
high-performance computing communities. Recent work on trace compression,
pattern mining, and database-backed profiling shares TraceLoom's goal of
reducing raw trace complexity. Our contribution is the specific adaptation to
LLM serving traces, where repeated decode iterations and model layers create
a distinctive hierarchical structure that can be exploited for compression and
attribution.

\textbf{Design-level profiling studies.} Prior profiling studies of data
stream processing systems~\cite{dsp-profiling} established the methodology of
abstracting common designs and studying their interaction with modern hardware
through mechanism-driven benchmarks. Our work adapts this methodology to LLM
serving, replacing DSP-specific designs (pipelined message passing, data
parallelism, JVM runtime) with serving-specific designs (KV state management,
token scheduling, hardware execution, program-aware serving).


% ============================================================
% 10. CONCLUSION
% ============================================================
\section{Conclusion}

Modern LLM serving frameworks should be profiled as online stateful runtime
systems, not as simple inference wrappers. By combining design-level workloads
with structured trace attribution, we can move beyond throughput rankings and
ask whether KV state, scheduling, communication, and program-level
optimizations actually change the intended execution structures.

We identified four recurring design principles---KV state management, online
token scheduling, hardware-aware execution, and program-aware serving---and
formulated mechanism-driven workloads as probes for each. Through TraceLoom,
we showed that raw NPU profiler timelines can be compressed into repeated
execution structures that expose communication neighborhoods, auxiliary
prelude costs, and cross-run optimization signals. Our perturbation case
studies demonstrated that source-plausible optimizations do not necessarily
produce stable trace-level changes, motivating trace-level optimization gates
as a complement to macro benchmark regression testing.

Future work should extend the methodology to additional serving frameworks
(SGLang, TensorRT-LLM), broader workload coverage (structured output,
multi-turn agents), and cross-platform trace adapters (NVIDIA Nsight/CUPTI).
As LLM applications grow in complexity, profiling must evolve to expose the
runtime abstractions---KV lifecycle, scheduler iterations, program
structure---that govern serving performance.


% ============================================================
% REFERENCES
% ============================================================
\begin{thebibliography}{00}

\bibitem{vllm}
W. Kwon et al., ``Efficient Memory Management for Large Language Model
Serving with PagedAttention,'' in \emph{Proc. SOSP}, 2023.

\bibitem{sglang}
L. Zheng et al., ``SGLang: Efficient Execution of Structured Language Model
Programs,'' in \emph{Proc. NeurIPS}, 2024.

\bibitem{tensorrtllm}
NVIDIA, ``TensorRT-LLM: A Framework for Optimizing Large Language Model
Inference,'' NVIDIA Developer Documentation, 2024.

\bibitem{tgi}
Hugging Face, ``Text Generation Inference,'' \url{https://github.com/huggingface/text-generation-inference}.

\bibitem{orca}
G. Yu et al., ``Orca: A Distributed Serving System for Transformer-Based
Generative Models,'' in \emph{Proc. OSDI}, 2022.

\bibitem{sarathi}
A. Agrawal et al., ``Sarathi: Efficient LLM Inference by Pinpointing
Inefficiencies with Communication-Aware Scheduling,'' in \emph{arXiv preprint}, 2023.

\bibitem{flashattn}
T. Dao et al., ``FlashAttention: Fast and Memory-Efficient Exact Attention
with IO-Awareness,'' in \emph{Proc. NeurIPS}, 2022.

\bibitem{burstgpt}
BurstGPT Project, ``BurstGPT: Failing to Forward Does Not Mean Falling
Behind,'' \url{https://github.com/HKU-DSG/BurstGPT}.

\bibitem{cann}
Huawei, ``CANN Profiling Tools,'' Ascend Computing Language documentation.

\bibitem{nsight}
NVIDIA, ``Nsight Systems User Guide,'' NVIDIA Developer Documentation.

\bibitem{perfetto}
Perfetto Project, ``Perfetto: System profiling, app tracing and trace
analysis,'' \url{https://perfetto.dev/}.

\bibitem{dsp-profiling}
A. Ilyas et al., ``Revisiting the Design of Data Stream Processing Systems
on Modern Multi-Core CPUs,'' in \emph{Proc. IEEE ICDE}, 2017.

\end{thebibliography}

\end{document}
